%%%%%%%%%%%%%%%%%%%%%%%%%%%%%
%%%%%%%%%%%%%%%%%%%%%%%%%%%%%

The Learning With Errors (LWE) problem was introduced by Oded Regev in 2005~\cite{STOC:Regev05} as a noisy generalization of the learning parity with noise problem over finite fields.


% -------------------------------------------
% OFFICIAL DEFINITION by LUCA
% -------------------------------------------
% \begin{definition}[LWE~\cite{EPRINT:AlbPlaSco15,STOC:Regev05}, Ring-LWE~\cite{EC:LyuPeiReg10} and Module-LWE~\cite{EPRINT:BraGenVai11}]
%     Let $N \in \Z^+$ (usually a power of $2$), $n, q \in \Z^+$, 
% $R_q = \Z_q[x]/\langle h_N(x) \rangle$, where $h_N(x)$ is an irreducible polynomial of degree $N$ (usually the $N$-th cyclotomic polynomial, e.g. $x^N + 1$), and let $\chi$ be an error distribution over $R_q$.
% For a secret vector $\vars{s} \in R_q^n$, the distribution
% $A_{\vars{s},\chi}$ produces samples $(\vars{a}, b) \in R_q^n \times R_q$ where:
% \[
%     \vars{a} \xleftarrow{\$} R_q^n, \quad e \leftarrow \chi, \quad
%     b = \langle \vars{a}, \vars{s} \rangle + e \in R_q.
% \]
% The \emph{search} variant asks to recover $\vars{s}$ given arbitrarily many samples from
% $A_{\vars{s},\chi}$. The \emph{decision} variant asks to distinguish
% $A_{\vars{s},\chi}$ from the uniform distribution over $R_q^n \times R_q$.

% If $n=1$, we refer to the problem as Ring-LWE (\rlwe), and if $n>1$, we refer to it as Module-LWE (\mlwe). By replacing \(R_q\) with \(\Z_q\), we can also express the standard LWE problem.
% \end{definition}

% -------------------------------------------
% Another trial 
% -------------------------------------------
\begin{definition}[LWE~\cite{EPRINT:AlbPlaSco15,STOC:Regev05}, Ring-LWE~\cite{EC:LyuPeiReg10} and Module-LWE~\cite{EPRINT:BraGenVai11}]
    Let $N, n, q \in \Z^+$, $R_q = \Z_q[x]/\langle h_N(x) \rangle$, where $h_N(x)$ is an irreducible polynomial of degree $N$, and let $\chi$ be an error distribution over $R_q$.
    For a secret vector $\vars{s} \in R_q^n$, the distribution $A_{\vars{s},\chi}$ produces samples $(\vars{a}, b) \in R_q^n \times R_q$ where:
    \[
        \vars{a} \xleftarrow{\$} R_q^n, \quad e \leftarrow \chi, \quad b = \langle \vars{a}, \vars{s} \rangle + e \in R_q.
    \]
    The \emph{search} variant asks to recover $\vars{s}$ given arbitrarily many samples from $A_{\vars{s},\chi}$. The \emph{decision} variant asks to distinguish $A_{\vars{s},\chi}$ from the uniform distribution over $R_q^n \times R_q$. In cryptographic applications, usually $h_N(x)=x^N+1$ where $N$ is a power of $2$.

    When $n=1$, we refer to the problem as Ring-LWE (\rlwe), and if $n>1$, we refer to it as Module-LWE (\mlwe). By replacing \(R_q\) with \(\Z_q\), we get the standard LWE problem.

    
\end{definition}


% % TODO: introduce notatino rlwe e mlwe
% \begin{definition}[LWE~\cite{EPRINT:AlbPlaSco15,STOC:Regev05}, Ring-LWE~\cite{EC:LyuPeiReg10}, Module-LWE~\cite{EPRINT:BraGenVai11}]
%     Let $N \in \Z_{>0}$ (usually a power of $2$), $q, n \in \Z_{>0}$, $R_q = \Z_q[x]/\langle h_N(x) \rangle$ with $h_N(x)$ an $N$-degree irreducible polynomial, and let $\chi$ be a probability distribution on $\Z$ (for LWE) or $R_q$ (for Ring-LWE and Module-LWE).
    
%     For a secret vector $\vars{s} \in \Z_q^n$ (for LWE) or $\vars{s} \in R_q^n$ (for Ring-LWE and Module-LWE), the distribution $A_{\vars{s},\chi}$ produces samples $(\vars{a}, b) \in R_q^n \times R_q$ (or $\Z_q^n \times \Z_q$ for LWE) where:
%     \[
%     \vars{a} \xleftarrow{\$} R_q^n, \quad e \leftarrow \chi, \quad
%     b = \langle \vars{a}, \vars{s} \rangle + e \in R_q.
%     \]
   

%     The \emph{search} variant asks to recover $\vars{s}$ given arbitrarily many samples from
%     $A_{\vars{s},\chi}$. The \emph{decision} variant asks to distinguish
%     $A_{\vars{s},\chi}$ from the uniform distribution over $R_q^n \times R_q$.

%     For $n=1$, Module-LWE (\mlwe) becomes Ring-LWE (\rlwe).
        
    
% \end{definition}



  
Since its introduction, the LWE problem is investigated in algorithm and cryptography literature. In particular, LWE with \emph{bounded errors} (LWBE) was investigated in~\cite{EPRINT:ACFP14,ICALP:AroGe11,EC:NMSU25,EPRINT:Steiner24b}. Introduced by Arora and Ge~\cite{ICALP:AroGe11}, an \emph{algebraic algorithm} generates a system of error-free multivariate polynomial equations from LWE samples and solve the system using linearization technique. In~\cite{EPRINT:ACFP14}, the complexity of Arora-Ge technique was improved by using Gr{\"o}bner basis methods to solve the polynomial system of equations. Recent analysis~\cite{EC:NMSU25,EPRINT:Steiner24b} have further improved the complexity of this approach.

The \emph{lattice-based algorithms}, 
using primal and dual approaches that reduce LWE to BDD/uSVP or SIS and rely on lattice reduction techniques such as BKZ and hybrid methods~\cite{EC:Albrecht17,EPRINT:AlbPlaSco15,RSA:LinPei11}, are the most efficient algorithms for solving LWE. The \emph{combinatorial algorithms}, notably BKW-based~\cite{DCC:ACFFP15} methods, reduce the problem dimension at the cost of increased noise and are mainly effective for instances with small secrets or large sample sizes. In~\cite{EPRINT:AlbPlaSco15}, the summary of a wide range of algorithms for LWE problem were given. 

%%%%%%%%%%%%%%%%%%%%%%%%%%%%%

\noindent
\textbf{Hardness and Cryptographic Application.}
The foundational hardness results such as reduction of worst-case approximate GapSVP problem to the average-case LWE problem~\cite{STOC:BLPRS13,DBLP:conf/stoc/Peikert09,STOC:Regev05,DBLP:journals/jacm/Regev09}, establish LWE as a central hardness assumption in lattice-based cryptography and a foundation for post-quantum cryptographic constructions.
% Regev proved that LWE is as hard on average as worst-case lattice problems such as GapSVP and SIVP via quantum reductions~\cite{STOC:Regev05}.
% The hardness guarantees was later extended to purely classical reductions by Brakerski et al.~\cite{STOC:BLPRS13}, under specific, but practically relevant, parameter choices. \ar{need to be specific}
%It has been shown that worst-case GapSVP reduces to average-case LWE~\cite{STOC:BLPRS13,STOC:Regev05}.
%%%%%%%%%%%%%%%%%%%%%%%%%%%%%
%The LWE problem has proven to be a rich and versatile source of many innovative cryptographic constructions. 
Examples of such constructions include secure public-key encryption schemes resilient to both chosen-plaintext~\cite{RSA:LinPei11,DBLP:conf/crypto/PeikertVW08,STOC:Regev05} and chosen-ciphertext attacks~\cite{DBLP:conf/eurocrypt/MicciancioP12,DBLP:conf/stoc/Peikert09,DBLP:conf/stoc/PeikertW08}, oblivious transfer protocols~\cite{DBLP:conf/crypto/PeikertVW08}, identity-based encryption~\cite{DBLP:conf/eurocrypt/AgrawalBB10,DBLP:conf/crypto/AgrawalBB10,DBLP:conf/eurocrypt/CashHKP10,DBLP:conf/stoc/GentryPV08}, leakage-resilient cryptography~\cite{DBLP:conf/tcc/AkaviaGV09,DBLP:conf/crypto/ApplebaumCPS09,DBLP:conf/innovations/GoldwasserKPV10}, fully homomorphic encryption~\cite{DBLP:conf/crypto/Brakerski12,DBLP:conf/innovations/BrakerskiGV12,DBLP:conf/focs/BrakerskiV11}, among many others.
Structured variants, including Ring-LWE~\cite{EC:LyuPeiReg10} and Module-LWE~\cite{EPRINT:BraGenVai11}, improve efficiency while retaining worst-case hardness, with MLWE underlying practical schemes such as CRYSTALS-Kyber~\cite{NISTPQC:CRYSTALS-KYBER22}.


\subsection{Our Result}

%% Table 1
% \begin{table}[hb!]
%     \centering
%     \renewcommand{\arraystretch}{1.7} 
%     \begin{tabular}{
%     >{\arraybackslash}m{1cm}
%     >{\arraybackslash}m{1cm}
%     >{\arraybackslash}m{1.5cm}
%     >{\arraybackslash}m{4cm}|
%     >{\arraybackslash}m{3.5cm}
%     >{\arraybackslash}m{3cm}
%     }
%     \hline
%     & &  \multicolumn{2}{c}{This work} & \multicolumn{2}{c}{Arora and Ge~\cite{ICALP:AroGe11}} \\
%     \hline
%   & Secret size & No. Samples & Time Complexity & No. Samples & Time Complexity\\ 
% \hline\hline
% LWE & $n$  & $\binom{n+d-1}{d}$ & $\bigO\left(n^{\max(3(d-1), d\omega)}\right)$ & $\bigO\left(q\log(q)n^d\right)$ & $\bigO\left(q\log(q)n^{d\omega}\right)$\\
% \hline
% RLWE & $N$  & $\frac{\binom{N+d-1}{d}}{N}$ & $\bigO\left(N^{\max(3(d-1), d\omega)}\right)$ & $\bigO\left(q\log(q)N^{d-1}\right)$ & $\bigO\left(q\log(q)N^{d\omega}\right)$\\
% \hline
% MLWE & $Nn$  & $\frac{\binom{Nn+d-1}{d}}{N}$ & $\bigO\left((Nn)^{\max(3(d-1), d\omega)}\right)$ & $\bigO\left(nq\log(q)(Nn)^{d-1}\right)$ & $\bigO\left(q\log(q)(Nn)^{d\omega}\right)$\\
% \hline\hline
%     \end{tabular}
%     \caption{Let $N, n, q \in \Z_{>0}$. The table shows the complexity of our algebraic attack on the three LWE variants with bounded errors (e.g. error entries/coefficients drawn from a set of size $d$). $\omega$ is the linear algebra constant ($2 \le \omega \le 3$) and the complexity is measured in the number of scalar multiplications in $\Z_q$.
%     }\label{tab:comparison}
% \end{table}
We consider algebraic algorithms for the search-LWE instances with \emph{bounded errors}, which asks to recover $\vars{s}$ from $(\vars{a}_i, b_i) = (\vars{a}_i, \langle\vars{a}_i, \vars{s}\rangle + e_i) \in \Z_q^n \times \Z_q$, for $1 \leq i \leq m$, where each $e_i$ lies in a fixed set of size $d$.
LWBE is a variant that remains as hard as standard LWE for suitable parameters~\cite{MicPei13}. The distinctness of our (algebraic) approach is its applicability
to general RLWE. While previous algebraic algorithms and their refinements uses polynomial system solving techniques that are applicable when the co-efficients are in $\F_q$. Surprisingly, previous works do not discuss the applicability of the techniques when the polynomial coefficients are in a ring, while it is well-known that it is non-trivial to extend the well-known polynomial system solving techniques such as Gr\"obner basis, when the coefficients are elements of a ring (e.g. $\Z_q$). 

\begin{table}[htbp]
    \centering
    \renewcommand{\arraystretch}{1.8}
    \resizebox{\textwidth}{!}{%
    \begin{tabular}{ c | c | c c | c c }
        \hline\hline
        \multirow{2}{*}{Variant} & \multirow{2}{*}{\makecell{Secret\\Size}} & \multicolumn{2}{c|}{This Work} & \multicolumn{2}{c}{Arora and Ge~\cite{ICALP:AroGe11}} \\
        \cline{3-6}
        & & No. Samples & Time Complexity & No. Samples & Time Complexity \\
        \hline\hline
        LWE  & $n$  & $\binom{n+d-1}{d}$           & $\bigO\left(n^{\max(3(d-1), d\omega)}\right)$      & $\bigO\left(q\log(q)n^d\right)$        & $\bigO\left(q\log(q)n^{d\omega}\right)$ \\
        \hline
        RLWE & $N$  & ${\binom{N+d-1}{d}}/{N}$  & $\bigO\left(N^{\max(3(d-1), d\omega)}\right)$      & $\bigO\left(q\log(q)N^{d-1}\right)$    & $\bigO\left(q\log(q)N^{d\omega}\right)$ \\
        \hline
        MLWE & $Nn$ & ${\binom{Nn+d-1}{d}}/{N}$ & $\bigO\left({(Nn)}^{\max(3(d-1), d\omega)}\right)$   & $\bigO\left(nq\log(q){(Nn)}^{d-1}\right)$ & $\bigO\left(q\log(q){(Nn)}^{d\omega}\right)$ \\
        \hline\hline
    \end{tabular}%
    }
    \vspace{0.2cm} 
    \caption{Let $N, n, q \in \Z_{>0}$. The table shows the complexity of our algebraic attack on the three LWE variants with bounded errors (e.g., error entries/coefficients drawn from a set of size $d$). Here, $\omega$ is the linear algebra constant ($2 \le \omega < 3$) and the complexity is measured in the number of scalar multiplications in $\Z_q$.}\label{tab:comparison}
\end{table}



In this work, we introduce a new algebraic method for solving the \slwe\ problem which differs fundamentally from existing approaches. Instead of directly applying Gr\"obner basis method,
we use \emph{a combination of diagonalization and $S$-polynomial generation} techniques. We carry out a direct complexity analysis of our new algorithm.

Unlike previous Gr{\"o}bner basis based analyses~\cite{EPRINT:ACFP14,EPRINT:Steiner24b}, our method does not rely on semi-regularity assumptions, or degree of regularity estimation for deriving the complexity bound. Thus, our algorithm and its complexity analysis is not affected by the non-existence of semi-regular system or difficulty of estimating degree of regularity (precisely). The only assumption we make is that the LWE samples are random, which allows us to rely on linear independence properties that hold with high probability. %In particular, we do not assume semi-regularity or generic behavior of the polynomial system.

Despite the absence of the above mentioned assumptions, the complexity of our algorithm shows polynomial improvement (\Cref{tab:comparison}) over existing algebraic algorithms based on Gr\"obner basis method. Note that the requirement of a large number of samples for our algorithm remains the same as in existing algebraic algorithms. A detailed comparison with prior algebraic attacks is provided in Table~\ref{tab:comparison}.
Our main result is summarized in \Cref{thm:result}:

\begin{theorem}[Main Result]\label{thm:result}
The complexity of solving a \slwe\ instance (for bounded errors) with parameters $n$, $d$ $(<n)$ and $m_{\text{LWE}} = \binom{n+d-1}{d}$ samples is $\mathcal{O}\!\left(n^{\max(3(d-1), d\omega)}\right)$ scalar multiplications in $\Z_q$, where $n$ denotes the length of the secret vector to be recovered, $d$ the size of the set from which the error is drawn and $\omega$ is the linear algebra constant ($2 \le \omega \le 3$).
\end{theorem}

\begin{corollary}[Algorithm's complexity for \rlwe\ and \mlwe]
    \label{cor:result-rlwe-mlwe}
The complexity of solving a \rlwe\ instance (for bounded errors) with parameters $N, d$ $(<N)$ and $m_{\rlwe} = \binom{N+d-1}{d}/N$ samples is $\mathcal{O}\!\left(N^{\max(3(d-1), d\omega)}\right)$ scalar multiplications in $\Z_q$. Similarly, the complexity of solving a \mlwe\ instance with parameters $N, n, d$ $(<N)$ and $m_{\mlwe} = \binom{Nn+d-1}{d}/N$ samples is $\mathcal{O}\!\left((Nn)^{\max(3(d-1), d\omega)}\right)$ scalar multiplications in $\Z_q$.
\end{corollary}

% In particular, with $m = \binom{n+d-1}{d}$ samples and error support of size $d$, our complexity bound improves upon existing algebraic algorithms~\cite{EPRINT:ACFP14,EC:NMSU25,EPRINT:Steiner24b,ACISP:SunTibAbe20} for bounded error LWE.\@


\subsection{Overview of Previous Algebraic Algorithms}
\label{subsec:previous-algorithms}

\noindent
\textbf{Arora and Ge Technique~\cite{ICALP:AroGe11}.} 
Each sample $(\vars{a}_i, b_i) = (\vars{a}_i, \langle\vars{a}_i, \vars{s}\rangle + e_i) \in \Z_q^n \times \Z_q$, for $1 \leq i \leq m$, allows to generate a non-linear equation of degree $d$ in the $n$ components of the secret $\vars{s}$. %which holds with probability $1 - e^{\mathcal{O}({-t^2})}$.
Considering that a common choice for the error distribution is the discrete Gaussian with mean $0$ and standard deviation $\sigma$, denoted as $\mathcal{N}(0,\sigma)$,
we have
\[
\mathbb{P}\big[ e \overset{\scriptscriptstyle\$}{\gets} \chi : |e| > t \cdot \sigma \big] \leq \frac{2}{t\sqrt{2\pi}} e^{-t^2/2}\in e^{\mathcal{O}({-t^2})}, \quad \textrm{for all } t>0.
\]
A sample from a Gaussian distribution take values on the interval $[-t\cdot\sigma,\dots,t\cdot\sigma]$ of $\Z_q$ with probability $1 - e^{\mathcal{O}({-t^2})}$ when representing elements in $\Z_q$ as integers in $[-\lfloor q/2 \rfloor, \dots, \lfloor q/2 \rfloor]$.
When $e \overset{\scriptscriptstyle \$}{\gets} \mathcal{N}(0,\sigma)$ then $f(e)=0$ for $f(x) = x \cdot \prod_{i=1}^{t\cdot\sigma} (x+i)(x-i)$ with probability at least $1 - e^{\mathcal{O}({-t^2})}$.
The degree of $f$ is $d := 2 \cdot t \sigma + 1 < q$. %We denote by $n$ the dimension of the secret. %and by $m$ the number of samples. %and $d := 2 \cdot t + 1$.
The resulting polynomial system $\mathcal{F}_{\mathrm{LWE}}$ is solved using the linearization method, namely by replacing monomials with a new linearized variable and then solving the resulting linear system of equations. %The \belwe\ where $\vars{e} \in \{0,1\}^m$ was introduced in~\cite{MicPei13}. 
The \belwe\, in which $\vars{e} \in \{0,1\}^m$, is a particular instance and was introduced in~\cite{MicPei13}.
%%%%%%%%%%%%%%%%%%%%%%%%%%%%%
%%%%%%%%%%%%%%%%%%%%%%%%%%%%%


\noindent
\textbf{Application of Gr{\"o}bner Basis Method.}
The Gr{\"o}bner basis algorithm was applied %F4~\cite{faugere:F4} and F5~\cite{faugere:F5} can be applied
to solve the (error free) polynomial system, in~\cite{EPRINT:ACFP14}, resulting in significant speed-up over the Arora-Ge algorithm. For providing a
complexity estimate \cite{EPRINT:ACFP14} and susequent works \cite{ACISP:SunTibAbe20, EPRINT:Steiner24b} derived bounds on the so-called degree of regularity under semi-regularity assumption. While more recent result ~\cite{EC:NMSU25} derived a bound that holds with high probability.  
%For applying Gr{\"o}bner algorithm, or its variants, one should work over $\Z_q$
% The \emph{degree of regularity} ($d_{reg}$) and semi-regularity plays an important role in the complexity analysis of the existing algebraic algorithms (for LWE) that
% uses Gr{\"o}bner basis method. 
% While computing $d_{reg}$ for a polynomial system is difficult, for semi-regular system one can estimate it.  
% Assuming that the polynomial system (from LWE) is semi-regular, the complexity analysis (of Gr{\"o}bner basis algorithm) in~\cite{EPRINT:ACFP14}, provides estimates of the degree of
% regularity. Under the same assumption,~\cite{ACISP:SunTibAbe20} estimates $d_{reg}$ using Hilbert polynomial. While most random systems of polynomials are semi-regular, Fr\"oberg
% showed~\cite{DBLP:journals/jsc/FrobergH94} that there are polynomial systems that can never be semi-regular, and~\cite{HODGES2017519} showed that (for certain parameter choices) semi-regular sequence
% can never exist.

% Understanding and analyzing the degree of regularity is crucial for deriving fine-grained complexity bounds for Gr\"obner basis computations, and has therefore been the focus of several works aiming at increasingly accurate complexity estimates.
%%%%%%%%%%%%%%%%%%%%%%%%%%%%%
%%%%%%%%%%%%%%%%%%%%%%%%%%%%%


% Recently,~\cite{EPRINT:Steiner24b} provided a refined complexity analysis using solving degree and assuming the degree of regularity to be same as the degree of polynomial. While this result does not make semi-regularity assumption, still relies on the assumption of degree of regularity. In~\cite{EC:NMSU25}, under semi-regularity assumption, the $d_{reg}$ was shown to be $d$ with high probability. 
% and consequently gives a refined complexity analysis of LWE (by taking into account key parameters such as the modulus $q$, noise distribution, dimension $n$, and number of samples $m$). 
% Using concepts of \emph{generic coordinates} and \emph{solving degree}, this work reduces reliance on the semi-regularity assumptions of earlier analyses~\cite{EPRINT:ACFP14}.
% %Despite these refinements, the degree of regularity remains high for standard cryptographic parameters, implying very large algebraic complexity.
% In particular, the main contribution relies on providing a tighter bound on the estimate of the complexity of computing Gr\"obner basis.

% We emphasize that our approach does not aim at improving existing bounds on the degree of regularity; instead, we avoid this notion altogether by providing a fine-grained complexity directly in terms of the LWE parameters.

All the above-mentioned algebraic analysis that use polynomial system solving methods, work for over $\F_q$. 
%%%%%%%%%%%%%%%%%%%%%%%%%%%%%
%%%%%%%%%%%%%%%%%%%%%%%%%%%%%


\paragraph{Algorithms for RLWE/PLWE (Summary in \ref{appendix:summary})}
Because an RLWE sample can be naturally represented as $n$ standard LWE samples, any known attack against LWE can be directly applied to the RLWE setting. However, the added algebraic structure of polynomial rings introduces unique attack vectors that do not exist in standard LWE.
A primary class of these algebraic attacks exploits ring homomorphisms. Initially, it was demonstrated that PLWE exhibits weak instances when the defining polynomial satisfies $f(1) \equiv 0 \pmod q$. By ``guessing'' the evaluation of the secret polynomial at 1, an attacker can solve the Decision PLWE problem in $\mathcal{O}(q)$ time. While this targets highly specific and non-standard instances, subsequent work generalized this vulnerability to the broader RLWE problem. Specifically, researchers showed that by choosing an alternative polynomial basis, apparently secure RLWE instances can be mapped into these weak PLWE instances. If the error distortion during this transformation is sufficiently small, the $\mathcal{O}(q)$ attack remains effective. 

Furthermore, while these initial techniques were limited to the Decision variant of the problem by mapping to simple prime fields $\mathbb{F}_q$, later cryptanalysis extended these attacks to the Search RLWE problem. By utilizing homomorphisms to higher-degree finite fields $\mathbb{F}_{q^f}$ and applying a chi-square statistical test to the mapped error distributions, attackers successfully exploited ``subfield vulnerabilities'' in prime cyclotomic rings. This advancement established a Search-to-Decision reduction, enabling full secret key recovery when the chosen error variance is small. 

In a broader context, the authors in \cite{...} proved a polynomial-time equivalence between the PLWE and RLWE problems across a wide class of number fields. This guarantees that an algebraic weakness discovered in one representation compromises the others. Consequently, to ensure the security of all variants, modern cryptographic implementations strictly avoid arbitrary or prime cyclotomics. Instead, they utilize 2-power cyclotomic rings (where these problems coincide and are immune to the $\mathcal{O}(q)$ and subfield vulnerabilities) alongside cryptographically large moduli to render such attacks unfeasible in practice.